\section{Introduction}

The long-only minimum variance portfolio allocates capital across $n$ assets to minimize
portfolio return variance subject to two constraints: the weights must sum to one (the
portfolio is fully invested) and must be non-negative. The only input required in theory
is the $n \times n$ covariance matrix $\Sigma$ of the $n \times 1$ vector of asset
returns. In practice, $\Sigma$ is unknown and must be estimated from historical data.
To this end, one uses a $T\times n$ matrix of demeaned returns $X$, where $T$ denotes
the length of the estimation window. The classical estimator of $\Sigma$ is the sample
covariance matrix $\hat \Sigma = X^\top X/T$.

The Lagrangian for the problem
\[
\min_w\; w^\top \hat \Sigma w \quad \text{subject to} \quad \mathbf{1}^\top w = 1,
\]
where $\mathbf{1}$ denotes a conformable vector of ones, is
\[
\mathcal{L}(w,\lambda) = w^\top \hat \Sigma w - \lambda(\mathbf{1}^\top w - 1),
\]
whose stationarity condition $2 \hat \Sigma w = \lambda\mathbf{1}$ immediately gives
$w \propto \hat \Sigma^{-1}\mathbf{1}$; this was made explicit by \cite{merton1972} in
the context of the efficient frontier. The budget constraint is recovered by the single
rescaling
\[
w = v/(\mathbf{1}^\top v), \qquad v = \hat \Sigma^{-1}\mathbf{1}.
\]
What has \emph{not} been exploited is this structure computationally: the standard
pipeline still forms $\hat \Sigma = X^\top X/T$ explicitly before solving.

This paper makes two contributions. First, we apply conjugate gradients (CG) to
$\hat \Sigma v = \mathbf{1}$ \emph{matrix-free}: $\hat \Sigma = X^\top X/T$ is never
assembled. Each CG iteration requires only two matrix-vector products with $X$, reducing
per-iteration cost from $O(n^2)$ to $O(nT)$ and eliminating $O(n^2)$ storage.
Non-negative weights are enforced by a primal-dual active-set loop that calls the CG
solver on successively smaller subproblems. Because CG is a Krylov method, it builds
successive iterates in the Krylov subspace $\mathcal{K}_k(\hat\Sigma, \mathbf{1})$ and
achieves $O(\sqrt\kappa)$ convergence, where $\kappa = \kappa(\hat\Sigma)$
is the spectral condition number (Proposition~\ref{prop:cg}). This $\sqrt\kappa$ rate is the defining advantage
of Krylov methods over general first-order schemes, which converge in $O(\kappa)$
iterations without exploiting the subspace structure (Proposition~\ref{prop:proximal}).

Second, we show that Ledoit-Wolf shrinkage~\cite{ledoit2004}, proposed to yield a more
accurate estimator of the true covariance matrix $\Sigma$ than the sample covariance
matrix, plays a dual role in this framework. Statistically, it reduces estimation error
by shrinking sample eigenvalues toward their mean. Numerically, the same modification
compresses the eigenvalue spread of the system matrix and directly accelerates CG
convergence (Proposition~\ref{prop:shrinkage}). On S\&P~500 equity data, LW shrinkage reduces the number of CG iterations
from 744 to 82. Statistical regularization and spectral preconditioning are, in this
problem, the same operation applied to the same matrix---a connection that does not
appear to have been made explicit in the existing literature.

CG of Hestenes and Stiefel~\cite{hestenes1952} is a cornerstone of large-scale
scientific computing; see, e.g., Golub and Van Loan~\cite{golub2013}. To the best of
our knowledge it has not been applied to portfolio optimisation in the matrix-free form
described here. The closest related work applies GMRES to the minimum variance problem
when $\hat \Sigma$ is singular~\cite{bajeux2012}; that formulation still assembles
$X^\top X$ and does not exploit the SPD structure or consider preconditioning.
General-purpose solvers such as Clarabel~\cite{clarabel2024} and OSQP~\cite{osqp} handle
the long-only constraint natively but require an assembled covariance matrix as input.
Markowitz's Critical Line Algorithm~\cite{markowitz1956} traces the full efficient
frontier via simplex-like pivots at $O(n^2)$ each; it targets the entire frontier
rather than a single minimum-variance solve and does not exploit the matrix-free
structure of $\hat\Sigma = X^\top X/T$.

For completeness, Section~\ref{sec:primaldual} also describes a simplex-projection
strategy based on proximal gradient descent. This method is \emph{not} a Krylov method:
it applies the gradient operator and projects onto the probability simplex at each step
without building a Krylov subspace, and consequently converges in $O(\kappa)$ rather than
$O(\sqrt\kappa)$ iterations. It is included as a matrix-free first-order comparator that
requires no outer active-set loop.

The remainder of the paper is organized as follows.
Sections~\ref{sec:problem}--\ref{sec:solvers} develop the problem formulation,
non-negativity strategies, Ledoit-Wolf preconditioning, and concrete solver
implementations. Section~\ref{sec:results} reports empirical results on S\&P~500
and synthetic data, benchmarking against direct KKT factorization, Clarabel, OSQP,
and proximal gradient.
